\subsection{Controlled-dimension and robustness supplements (E4-C, E5)}
\label{sec:supplements}

The E4-H supplement changed the design dimension, the spectral
parametrisation and the number of $R/T$ observables at the same time, so its
censored PASS could not be attributed to dimension alone. E4-C removes that
ambiguity. Both arms use the same RBF basis and width, the same true spectrum,
the same start point, the same 44 normalised observables, the same loss, the
same optimiser settings and the same evaluation budgets; the only factor that
differs is the number of free design variables. The 8-D spectral coefficients
are mapped into the 14-dimensional RBF coefficient space by a fixed
row-normalised Gaussian matrix, so the 8-D arm is an exact linear subspace of
the 16-D arm rather than a differently parametrised problem. The two synthetic
targets agree to $|\Delta y|_{\max} = 0$,
and a constant coefficient vector maps to itself to $1.1\times10^{-16}$.
The experiment is deliberately noise-free, because its purpose is to isolate
one factor; noise and multiple starts are handled separately in E5.

\begin{table}[t]
\centering
\caption{E4-C evaluation counts to reach a target loss. Each cell gives the
L-BFGS count over the Nelder--Mead count. Entries marked $\ge$ are censored
lower bounds, not measured ratios: the derivative-free run stopped at its
budget before reaching that loss.}
\label{tab:e4c}
\begin{tabular}{lll}
\hline
Target loss & 8-D & 16-D \\
\hline
0.01 & 1/1 = 1.00$\times$ (measured) & 1/1 = 1.00$\times$ (measured) \\
$10^{-3}$ & 7/25 = 3.57$\times$ (measured) & 7/60 = 8.57$\times$ (measured) \\
$10^{-4}$ & 25/---, $\ge$ 12.0$\times$ (censored) & 23/---, $\ge$ 13.0$\times$ (censored) \\
$10^{-5}$ & not reached & 33/---, $\ge$ 9.1$\times$ (censored) \\
$10^{-6}$ & not reached & 43/---, $\ge$ 7.0$\times$ (censored) \\
\hline
\end{tabular}
\end{table}

L-BFGS used 67 (8-D) and 64 (16-D)
forward/adjoint evaluations, and Nelder--Mead used 301 and 300
function calls, reaching best losses of $3.0\times10^{-4}$ and $2.2\times10^{-4}$.
Of the 10 arm--threshold combinations, 4
yield a measured ratio (2 of which are the trivial
first-call row at the loosest threshold) and 4 are censored.
Measured, uncensored comparisons are available at 1 informative threshold: at $10^{-3}$ the measured advantage is 3.57$\times$ (8-D) against 8.57$\times$ (16-D). The 16-D advantage is larger at every such threshold, consistent with a dimension effect under this fixed-observation protocol. At $10^{-3}$ L-BFGS used 7 and 7 calls in the two arms, whereas Nelder--Mead used 25 and 60. The growth of the ratio is therefore attributable to the derivative-free simplex degrading with dimension, not to an improvement of the gradient-based search. Because the 8-D arm is a subspace of the 16-D one, the higher-dimensional inverse problem is also less constrained, which contributes to the lower final loss ($1.5\times10^{-7}$ versus $2.7\times10^{-5}$). This is a single-start, noise-free, single-model result, not a general efficiency claim. At $10^{-4}$ and stricter losses the derivative-free run is censored in both arms, so the recorded 12.04$\times$ (8-D) and 13.04$\times$ (16-D) are lower bounds rather than measured ratios; the E4-H PASS at that threshold therefore remains a censored result.

E5 addresses a different question: whether the picture survives changes of
start point, observation noise and the choice of derivative-free comparator.
Two deterministically perturbed starts were run against 5
baselines (Nelder--Mead, Powell, differential evolution, dual annealing and an
in-house RBF--UCB surrogate), each with a budget of 80 true
function calls and 1\% relative Gaussian on synthetic observables on the synthetic observables. At the
common budget the smallest mean best loss was 6.540e-04 (Powell) and the largest was 1.001e-03 (DualAnnealing);
7 of 10 runs reached $10^{-3}$ and 0 of 10
reached $10^{-4}$. The noisy objective has a loss floor set by the noise
variance, so these runs cannot be used as evaluation-count comparisons against
the noise-free gradient-based results, and no speedup is claimed from them.
The RBF--UCB baseline is a lightweight in-house implementation, not a tuned
Bayesian-optimisation library.

Both supplements are independent additions. The pre-registered E4 verdict
(the $\ge10\times$ target is \textbf{not met} in eight dimensions) and the
E4-H censored lower-bound PASS are unchanged, as are the statements that all
targets are synthetic inverse crimes and that no SCAPS-1D cross-validation or
experimental calibration has been performed.
